\documentclass{standalone}
\usepackage{circuitikz}
\usetikzlibrary{calc}
\definecolor{circuitnotemark}{HTML}{FE00FE}
\begin{document}
\begin{circuitikz}[american, line width=0.8pt]
\ctikzset{bipoles/length=1.2cm}
\ctikzset{grounds/scale=1.36}
\coordinate (b1) at (0,-2);
\coordinate (e1) at (0,-8);
\coordinate (b3) at (4,-2);
\coordinate (b6) at (10,-2);
\coordinate (b7) at (12,-2);
\coordinate (e7) at (12,-8);
\coordinate (e8) at (14,-8);
\coordinate (b8) at (14,-2);
\coordinate (b10) at (18,-2);
\draw (b1) to[esource, l_=$\mathrm{W1}$] (e1); % line 3
\draw (-0.18,-5) -- ++(0,0.18) -- ++(0.18,0) -- ++(0,-0.36) -- ++(0.18,0) -- ++(0,0.18); % line 3
\node[ground] at (e1) {}; % line 4
\draw (b3) to[R, l_=$R_{1}$, a^=$1\,\mathrm{k}\Omega$] (b6); % line 5
\draw (b7) to[D, l_=$D_{1}$] (e7); % line 6
\draw (e8) to[D, l_=$D_{2}$] (b8); % line 7
\node[ground] at (e8) {}; % line 8
\draw (b10) node[ocirc]{} node[above left]{OUT}; % line 9
\draw (b1) -- (b3); % line 11
\draw (b6) -- (b7); % line 12
\draw (b7) -- (b8); % line 13
\draw (e7) -- (e8); % line 14
\draw (b8) -- (b10); % line 15
\node[circ] at (b7) {};
\node[circ] at (e8) {};
\node[circ] at (b8) {};
\node[anchor=south west, inner sep=2pt, yshift=4pt, circuitnotemark, font=\LARGE\bfseries] (circuittitle) at (current bounding box.north west) {X};
\path (circuittitle.south west) rectangle ++(6.218,0.853);
\node[anchor=north east, inner sep=2pt, circuitnotemark, font=\large] (circuitstamp) at ([yshift=-0.593cm]current bounding box.south east) {X};
\path (circuitstamp.north east) rectangle ++(-5.08,-0.593);
\end{circuitikz}
\end{document}
